\chapter{Исходный код разработанной системы}\label{appendix1}							% Заголовок
%\addcontentsline{toc}{chapter}{Second call for chapters to participate in the book Machine learning in analysis of biomedical and socio-economic data}	% Добавляем его в оглавление

\begin{lstlisting}[language=SQL,
	caption={SQL-запрос для анализа таблиц},
	label={lst:tables_info}]
	SELECT 
	n.nspname AS schema, 
	c.relname AS table_name, 
	c.relkind, 
	pg_size_pretty(pg_total_relation_size(c.oid)) AS table_size, 
	COALESCE(s.seq_scan, 0) AS seq_scans, 
	COALESCE(s.idx_scan, 0) AS idx_scans, 
	COALESCE(s.n_tup_ins, 0) AS inserts, 
	COALESCE(s.n_tup_upd, 0) AS updates, 
	COALESCE(s.n_tup_del, 0) AS deletes, 
	s.last_vacuum, 
	s.last_autovacuum, 
	s.last_analyze, 
	s.last_autoanalyze 
	FROM pg_class c 
	JOIN pg_namespace n ON n.oid = c.relnamespace 
	LEFT JOIN pg_stat_user_tables s ON s.relid = c.oid 
	WHERE c.relkind = 'r' 
	AND n.nspname NOT IN ('pg_catalog', 'information_schema') 
	ORDER BY pg_total_relation_size(c.oid) DESC; 
\end{lstlisting}

\begin{lstlisting}[language=bash,
	caption={Файл docker-compose.yml},
	label={lst:dockercompose}]
	version: '3.9'
	
	services:
	
	prometheus:
		image: prom/prometheus:latest
		container_name: monitoring-prometheus
		volumes:
			- ./prometheus/prometheus.yml:/etc/prometheus/prometheus.yml
			- ./prometheus/rules:/etc/prometheus/rules
		command:
			- --config.file=/etc/prometheus/prometheus.yml
		ports:
			- "${PROMETHEUS_PORT}:9090"
		networks:
			- monitoring-network
		restart: unless-stopped
	
	alertmanager:
		image: prom/alertmanager:latest
		container_name: monitoring-alertmanager
		volumes:
			- ./alertmanager/alertmanager.yml:/etc/alertmanager/alertmanager.yml
			- ./alertmanager/template.tmpl:/etc/alertmanager/template.tmpl
		command:
			- --config.file=/etc/alertmanager/alertmanager.yml
		ports:
			- "${ALERTMANAGER_PORT}:9093"
		networks:
			- monitoring-network
		restart: unless-stopped
	
	grafana:
		image: grafana/grafana:latest
		container_name: monitoring-grafana
		environment:
			- GF_SECURITY_ADMIN_USER=${GRAFANA_ADMIN_USER}
			- GF_SECURITY_ADMIN_PASSWORD=${GRAFANA_ADMIN_PASSWORD}
		volumes:
			- ./grafana/provisioning:/etc/grafana/provisioning
			- ./grafana/dashboards:/var/lib/grafana/dashboards
		ports:
			- "${GRAFANA_PORT}:3000"
		networks:
			- monitoring-network
		depends_on:
			- prometheus
		restart: unless-stopped
	
	node-exporter:
		image: prom/node-exporter:latest
		container_name: monitoring-node-exporter
		ports:
			- "9100:9100"
		networks:
			- monitoring-network
		restart: unless-stopped
	
	postgres-exporter:
		image: prometheuscommunity/postgres-exporter:master
		container_name: monitoring-postgres-exporter
		environment:
			- DATA_SOURCE_NAME=postgres://${GP_USER}:${GP_PASSWORD}@${GP_HOST}:${GP_PORT}/${GP_DATABASE}?sslmode=disable
		ports:
			- "9187:9187"
		networks:
			- monitoring-network
		extra_hosts:
			- "host.docker.internal:host-gateway"
		restart: unless-stopped
	
	greenplum-exporter:
		image: qinyugui/greenplum-exporter:1.1
		container_name: monitoring-greenplum-exporter
		environment:
			- GPDB_DATA_SOURCE_URL=postgres://${GP_USER}:${GP_PASSWORD}@${GP_HOST}:${GP_PORT}/${GP_DATABASE}?sslmode=disable
		ports:
			- "9188:9297"
		networks:
			- monitoring-network
		extra_hosts:
			- "host.docker.internal:host-gateway"
		restart: unless-stopped
	
	mkdocs:
		image: polinux/mkdocs:arm64v8-1.5.2
		container_name: monitoring-mkdocs
		volumes:
			- ./runbooks:/docs
		ports:
			- "8000:8000"
		networks:
			- monitoring-network
		restart: unless-stopped
	
	networks:
		monitoring-network:
	driver: bridge
\end{lstlisting}

\begin{lstlisting}[language=bash,
	caption={Единый конфигурационный файл .env.example},
	label={lst:envexample}]
	GP_HOST=
	GP_PORT=
	GP_USER=
	GP_PASSWORD=
	GP_DATABASE=
	
	EXPORTER_PORT=9187
	PROMETHEUS_PORT=9090
	ALERTMANAGER_PORT=9093
	GRAFANA_PORT=3000
	
	GRAFANA_ADMIN_USER=
	GRAFANA_ADMIN_PASSWORD=
	
	ALERT_EMAIL_TO=
	ALERT_EMAIL_FROM=
	
	SMTP_HOST=smtp.gmail.com
	SMTP_PORT=587
	SMTP_USER=
	SMTP_PASSWORD=
\end{lstlisting}

\begin{lstlisting}[language=bash,
	caption={Файл generate-configs.sh},
	label={lst:generateconfigs}]
	#!/bin/bash
	
	set -e
	
	echo "Generating configs..."
	
	# Prometheus
	envsubst < prometheus/prometheus.yml.template > prometheus/prometheus.yml
	
	# Alertmanager
	envsubst < alertmanager/alertmanager.yml.template > alertmanager/alertmanager.yml
	
	# Grafana datasource
	envsubst < grafana/provisioning/datasources/datasource.yml.template > grafana/provisioning/datasources/datasource.yml
	
	echo "Configs generated"
\end{lstlisting}


%

%






%% В случае, когда таблица (рисунок) размещаются на последней странице, для переноса названия приложения на новую строку используем:
%\NewPage % начать новое приложение с новой страницы 